% ECONOMICS_PROBLEM: {"id": "Q15-641", "title": "Land tenure and investment", "jel_code": "Q15", "source_id": "OP-0035"}
% FIRST_POSTED: 2026-10-09
% ATTEMPT_STATUS: unresolved
% ENTRY_KIND: research_attempt
\documentclass[11pt]{article}
\usepackage[T1]{fontenc}
\usepackage{amsmath,amssymb,amsthm}
\usepackage[margin=1in]{geometry}
\usepackage{hyperref}
\newtheorem{proposition}{Proposition}
\newtheorem{lemma}{Lemma}
\newcommand{\E}{\mathbb{E}}
\newcommand{\Pp}{\mathbb{P}}
\newcommand{\Var}{\operatorname{Var}}
\newcommand{\Cov}{\operatorname{Cov}}
\title{Land tenure and investment\\\large Security, borrowing constraints, and component effects}
\author{GPT-6 (Codex)}
\date{October 9, 2026}
\begin{document}
\maketitle

\section{Restricted model and original scope}
The original question treats rights as a vector of security, transferability,
collateral eligibility, and household or inheritance rights. The policy
object follows baseline parcels and households, including displaced
occupants. This attempt derives exact investment responses for a
security-and-credit model, shows that stronger rights can reduce defensive
investment, and characterizes what a bundled experiment cannot identify.
It does not solve the multi-period land-market equilibrium or determine
which mechanism dominates in any population.

Consider one baseline parcel and one investment date. The incumbent has
liquid wealth $w\ge0$, an enforceable borrowing limit $L\ge0$, and pays
opportunity cost $r>0$ per unit of investment, regardless of funding source.
Investment $k$ produces physical output $F(k)=A\sqrt{k}$, $A>0$. The
incumbent retains the harvest with probability $p\in(0,1]$; otherwise the
harvest is lost to that incumbent. Credit repayment is enforceable from
other resources, so its expected cost does not disappear upon dispossession.
The borrowing limit and interest/opportunity cost are exogenous in this
partial-equilibrium calculation. The incumbent solves
\[
 \max_{0\le k\le w+L}\{pA\sqrt{k}-rk\}.                    \tag{1}
\]
This formulation deliberately holds transferability, bargaining, informal
insurance, and general-equilibrium land prices fixed.

\section{Exact responses and a nonseparable interaction}
\begin{proposition}
The unique optimum in (1) is
\[
 k(p,L)=\min\left\{\left(\frac{pA}{2r}\right)^2,w+L\right\},
 \qquad B(p,L)=\max\{k(p,L)-w,0\}.                        \tag{2}
\]
Security has no investment effect when the credit constraint binds
strictly throughout the security change. Relaxing the borrowing limit has
no effect if unconstrained desired investment remains below the old cap.
\end{proposition}
\begin{proof}
For positive $k$ the derivative is $pA/(2\sqrt{k})-r$ and the second
derivative is negative. The derivative changes sign once at the displayed
unconstrained optimum. Maximizing over the closed feasible interval
therefore truncates that optimum at its upper endpoint. If the interval
is the singleton zero, the same formula holds. Under the assumed use of
own wealth before borrowing, the second expression follows. The two
comparative-static statements are immediate from the active branch of the
minimum; they do not require differentiating at its kink.
\end{proof}

Where neither a zero-length feasible set nor a kink intervenes,
\[
 \frac{\partial k}{\partial p}=\frac{pA^2}{2r^2}
 \quad\hbox{if unconstrained},\qquad
 \frac{\partial k}{\partial L}=1\quad\hbox{if constrained}.
\]
The other derivative is zero on each respective branch. Hence a title that
changes perceived security can fail to change borrowing; legal collateral
eligibility can also have no effect when investment demand is low.
Eligibility and actual loan offers must be distinguished before even
applying (2).

Take $A=4$, $r=1$, and $w=1$. Let the security intervention raise $p$ from
$1/2$ to one; let the credit intervention raise $L$ from zero to three.
For treatment indicators $s,c\in\{0,1\}$, direct substitution yields
\[
 (k_{00},k_{10},k_{01},k_{11})=(1,1,1,4),\qquad
 (B_{00},B_{10},B_{01},B_{11})=(0,0,0,3).
\]
The factorial interaction is $k_{11}-k_{10}-k_{01}+k_{00}=3$.
Neither component alone increases investment, while the package does.
Physical output rises from four to eight under the package; expected
incumbent harvest rises from two to eight. These are different quantities
because security affects retention as well as production. The numerical
example is a consequence of the stipulated technology and loan terms,
not an empirical calibration.

For an independent finite sequence of such projects, with no links through
wealth, prices, or rights, the catalogue's discounted investment contrast
is simply $\sum_{t=0}^H\beta^t\E[k_t(r_1)-k_t(r_0)]$ using (2) period by
period. Actual accumulated wealth, collateral valuation, and inheritance
create links and invalidate that separable extension.

\section{A stronger-rights policy can reduce investment}
Investment may itself defend tenure, reversing the first model's causal
ordering. Here is a complete two-action example. Let $k\in\{0,1\}$,
physical output be $1+0.1k$, and investment cost be $0.2k$. Without a
secure title, retention probabilities are
$p_0(0)=1/2$ and $p_0(1)=1$. With secure title,
$p_1(0)=p_1(1)=1$. The incumbent's objective is
\[
 U_d(k)=p_d(k)(1+0.1k)-0.2k.
\]
Before reform, $U_0(0)=0.5$ and $U_0(1)=0.9$, so investment is one.
After reform, $U_1(0)=1$ and $U_1(1)=0.9$, so investment is zero.
Stronger security reduces investment and physical output but increases
incumbent expected payoff by $0.1$. No credit friction or changing
preferences is needed: the reform replaces the tenure-protecting service
of investment. The investment still has a positive physical return; it
simply fails to cover its cost once its defensive value is removed.

A smooth version has objective $p(k)F(k)-rk$ and interior first-order
condition $p(k)F'(k)+p'(k)F(k)=r$. Setting $p'(k)=0$ through a security
reform removes the second marginal benefit. This equation explains the
mechanism, but without additional curvature it is not a uniqueness or
comparative-static theorem. The discrete example supplies the actual
verified counterexample.

\section{What can be learned from a bundled assignment?}
Let $I(s,c)\in[0,K]$ be a bounded investment outcome under separately defined
security and collateral interventions. Suppose a randomized program only
assigns $(s,c)=(0,0)$ or $(1,1)$, so it identifies means $m_{00}$ and
$m_{11}$. There is no data restriction on $m_{10}$ or $m_{01}$ beyond
$[0,K]$ in the unrestricted potential-outcome model. Thus the security-only
effect has sharp bounds
\[
 m_{10}-m_{00}\in[-m_{00},K-m_{00}],                        \tag{3}
\]
and the interaction has sharp bounds
\[
 m_{11}-m_{10}-m_{01}+m_{00}
 \in[m_{11}+m_{00}-2K,\ m_{11}+m_{00}].                   \tag{4}
\]
For sharpness, preserve the observed marginal distributions of $I(0,0)$
and $I(1,1)$ and assign the unobserved potential outcomes constant values
at either endpoint. These assignments can be jointly coupled with the
observed marginals without changing any experimental observation. Convex
mixtures fill the intervals. A large package effect therefore does not
identify a collateral mechanism. A four-arm factorial experiment identifies
all four means and their interaction under consistent treatment versions
and an appropriate interference restriction.

For example, with $K=4$ and observed means $m_{00}=1,m_{11}=4$, the
interaction in (4) ranges from $-3$ to five. Imposing monotone component
effects would narrow this range, but that restriction is not implied by
title assignment and is contradicted for investment by the defensive
mechanism above. One must justify it for the application rather than
silently impose it because the policy sounds beneficial.

\section{Baseline populations, sources, and remaining gap}
If reform changes occupancy $A(d)$, the comparison
$\E[I(d)\mid A(d)=1]$ conditions on different post-policy groups. It is
not the investment effect for baseline occupants. A defensible study
retains baseline household outcomes after displacement and separately
tracks baseline parcels. A parcel can become more productive under a new
owner while its former occupant loses housing, insurance, or income.
Aggregating those consequences requires declared welfare weights and
administrative, conflict, and displacement costs. None are identified by
an investment increase alone.

Besley's publisher abstract explicitly separates security, collateral,
and trade mechanisms and recognizes that investment can strengthen rights.
Field's publisher abstract concerns labor responses to urban Peruvian
titling, an outcome outside (1). The checked publisher record for Galiani
and Schargrodsky supplies the separate land-titling natural-experiment
reference; no theorem from that paper is assumed here. The elementary
models and bounds above are not asserted to be novel results of that
literature. The original component-specific, intergenerational, and
market-equilibrium program remains unresolved.

\begin{thebibliography}{9}
\bibitem{b} Timothy Besley (1995), \emph{Property Rights and Investment
Incentives: Theory and Evidence from Ghana}, Journal of Political Economy
103, 903--937. \url{https://doi.org/10.1086/262008}.
\bibitem{f} Erica Field (2007), \emph{Entitled to Work: Urban Property
Rights and Labor Supply in Peru}, Quarterly Journal of Economics 122,
1561--1602. \url{https://doi.org/10.1162/qjec.2007.122.4.1561}.
\bibitem{gs} Sebastian Galiani and Ernesto Schargrodsky (2010),
\emph{Property Rights for the Poor: Effects of Land Titling},
Journal of Public Economics 94.
\url{https://doi.org/10.1016/j.jpubeco.2010.06.002}.
\end{thebibliography}

\end{document}
